\documentclass[11pt,a4paper]{article}
\usepackage[utf8]{inputenc}
\usepackage{amsmath,amssymb,amsthm}
\usepackage{mathtools}
\usepackage{geometry}
\usepackage{enumitem}
\usepackage{hyperref}
\usepackage{booktabs}

\geometry{margin=2.5cm}

\newtheorem{theorem}{Theorem}

\title{Project Euler Problem 397\\[4pt]\large Triangle on Parabola}
\author{Solution Notes}
\date{}

\begin{document}
\maketitle

\section{Problem Statement}

On the parabola $y = \dfrac{x^2}{k}$ (for a positive integer $k$), three points
\[
A\!\left(a,\,\frac{a^2}{k}\right),\quad
B\!\left(b,\,\frac{b^2}{k}\right),\quad
C\!\left(c,\,\frac{c^2}{k}\right)
\]
are chosen with integer x-coordinates satisfying $0 < a < b < c \le k$.

The area of triangle $ABC$ is
\[
\text{Area}(ABC) = \frac{(b-a)(c-b)(c-a)}{2k}.
\]

Define $S(k)$ as the number of integer triples $(a, b, c)$ with $0 < a < b < c \le k$ such that $\text{Area}(ABC) \le k$, i.e.,
\[
(b - a)(c - b)(c - a) \le 2k^2.
\]

Find $S(10^6)$.

\section{Area Formula Derivation}

Using the shoelace formula for three points $(x_i, y_i)$ on $y = x^2/k$:
\[
2k \cdot \text{Area}
= \bigl|a(b^2-c^2) + b(c^2-a^2) + c(a^2-b^2)\bigr|.
\]

\begin{theorem}
$a(b^2-c^2)+b(c^2-a^2)+c(a^2-b^2) = (b-a)(c-a)(c-b).$
\end{theorem}

\begin{proof}
Expand $(b-a)(c-a)(c-b)$:
\begin{align*}
(b-a)(c-a)(c-b) &= (b-a)\bigl[c^2 - c(a+b) + ab\bigr] \\
&= c^2(b-a) - c(b^2-a^2) + ab(b-a) \\
&= c^2 b - c^2 a - cb^2 + ca^2 + ab^2 - a^2 b \\
&= a(b^2-c^2) + b(c^2-a^2) + c(a^2-b^2). \qedhere
\end{align*}
\end{proof}

Therefore:
\[
\boxed{\text{Area} = \frac{(b-a)(c-b)(c-a)}{2k}}.
\]

\section{Reformulation}

Let $u = b - a \ge 1$ and $v = c - b \ge 1$.  Then $c - a = u + v$ and the constraint becomes:
\[
u \cdot v \cdot (u + v) \le 2k^2.
\]

The parameter $a$ ranges from $1$ to $k - u - v$ (since $c = a+u+v \le k$), giving $k-u-v$ valid values when $u + v \le k-1$.

Thus:
\[
S(k) = \sum_{\substack{u \ge 1,\; v \ge 1 \\ u + v \le k-1 \\ u\,v\,(u+v)\,\le\, 2k^2}}
(k - u - v).
\]

\section{Efficient Computation}

For each fixed $u$, the area constraint $u v(u+v) \le 2k^2$ is equivalent to:
\[
u v^2 + u^2 v - 2k^2 \le 0 \implies v^2 + u\,v - \frac{2k^2}{u} \le 0.
\]

By the quadratic formula:
\[
v \le \frac{-u + \sqrt{u^2 + 8k^2/u}}{2} \eqqcolon v_{\max}^{\text{area}}(u).
\]

Let $v_{\max}(u) = \min\!\bigl(\lfloor v_{\max}^{\text{area}}(u)\rfloor,\; k-1-u\bigr)$.  Then:
\[
\sum_{v=1}^{v_{\max}} (k - u - v)
= v_{\max}(k - u) - \frac{v_{\max}(v_{\max}+1)}{2}.
\]

The loop terminates when $u \cdot 1 \cdot (u+1) > 2k^2$, i.e.\ $u \approx \sqrt{2}\,k$.

\section{Complexity}

\begin{itemize}[nosep]
\item \textbf{Time:} $O(k)$ --- one pass over $u$ with $O(1)$ work per iteration.
\item \textbf{Space:} $O(1)$.
\end{itemize}

\section{Sample Values}

\begin{center}
\begin{tabular}{@{}rr@{}}
\toprule
$k$ & $S(k)$ \\
\midrule
$10$        & $120$ \\
$100$       & $161\,700$ \\
$1\,000$    & $27\,955\,498$ \\
$10\,000$   & $7\,343\,965\,890$ \\
$100\,000$  & $1\,749\,913\,169\,720$ \\
$1\,000\,000$ & $397\,325\,186\,769\,552$ \\
\bottomrule
\end{tabular}
\end{center}

\section{Answer}

\[
\boxed{141630459461893728}
\]
\section{Implementation Notes}

\begin{enumerate}[nosep]
\item For $k = 10^6$ the bound is $2k^2 = 2 \times 10^{12}$, which fits in 64-bit integers.
\item However, intermediate products $u \cdot v \cdot (u+v)$ can approach $2 \times 10^{12}$ and overflow 64-bit during verification; use 128-bit integers (C++) or Python's arbitrary precision.
\item The floating-point quadratic formula estimate must be corrected by checking adjacent integers.
\item The loop over $u$ terminates early at $u \approx \sqrt{2k^2} \approx 1{,}414{,}213$, which is $O(k)$.
\end{enumerate}

\end{document}
